\begin{table}[htbp]
\centering
\caption{ساختار و محتوای مداخله آموزشی در دو گروه پژوهش}
\label{tab:intervention}
{\small
\setlength{\tabcolsep}{3pt}
\renewcommand{\arraystretch}{1.2}
\begin{tabularx}{0.985\textwidth}{p{0.10\textwidth} X p{0.18\textwidth} p{0.22\textwidth}}
\hline
\textbf{مرحله} & \textbf{محتوای آموزشی} & \textbf{گروه آموزش کارگاهی} & \textbf{گروه کارگاهی تلفیق‌شده با خرده‌یادگیری} \\
\hline
جلسه اول & آشنایی با پرفشاری خون و عوارض آن، مفهوم خودمراقبتی و روش صحیح اندازه‌گیری فشار خون در منزل & به صورت حضوری (۶۰ دقیقه‌) & به صورت حضوری (۶۰ دقیقه‌) همراه با میکروویدیوهای ۱ و ۲ \\
جلسه دوم & تغذیه سالم، کاهش نمک و چربی، آشنایی با الگوی غذایی «\lr{DASH}» و کنترل وزن & به صورت حضوری (۶۰ دقیقه‌) & به صورت حضوری (۶۰ دقیقه‌) همراه با میکروویدیوهای 3 و 4 \\
جلسه سوم & فعالیت بدنی منظم، ورزش متناسب با توان فردی، مدیریت تنش و راهبردهای آرام‌سازی & به صورت حضوری (۶۰ دقیقه‌) &به صورت حضوری (۶۰ دقیقه‌) همراه با میکروویدیوهای 5 و 6 \\
جلسه چهارم & داروهای رایج پرفشاری خون، پایبندی دارویی، مدیریت خودمراقبتی در زندگی روزمره و جمع‌بندی & به صورت حضوری (۶۰ دقیقه‌) & به صورت حضوری (۶۰ دقیقه‌) همراه با میکروویدیوهای 7 و 8 \\
\hline
\end{tabularx}
}
\end{table}
